% Test-case specification TC-0011 (CTest test TC-0011-errobj).
\documentclass[letterpaper]{article}

\usepackage[letterpaper,margin=1in]{geometry}
\emergencystretch=3em

% Compile from the repository root so this path resolves to
% the pinned WSP submodule.
\input{wsp/testing/test-case-library.tex}

\begin{document}

\def\TCID{TC-0011}
\def\TCTitle{%
  The Err object}
\def\TCPurpose{%
  Verify the members of the \texttt{Err} object, through the interface as
  compiled code calls them (values returned directly) and through late
  binding.}
\def\TCPriority{\TCPriorityHigh}

\def\TCRequirementRef{%
  REQ-0004}
\def\TCDesignRef{%
  Not applicable.}
\def\TCFeatureRef{%
  \texttt{tests/\allowbreak{}vbrun\_\allowbreak{}errobj\_\allowbreak{}tests.\allowbreak{}cpp},
  \texttt{src/\allowbreak{}helpers\_\allowbreak{}errobj.\allowbreak{}cpp}}
\def\TCTestTechnique{%
  Requirements-based testing; both calling paths compared.}

\def\TCPreconditions{%
  The library and the test program are built
  (\texttt{scripts/\allowbreak{}Invoke-Build.\allowbreak{}ps1}) and the test
  is registered with CTest.}
\def\TCEnvironment{%
  Windows with the toolchain of
  \texttt{scripts/\allowbreak{}Invoke-Build.\allowbreak{}ps1}; the local
  gate builds and runs x64 Debug, CI builds and runs the other
  architectures. The library is built from this repository in
  \texttt{out/\allowbreak{}build}.}
\def\TCAssumptions{%
  Not applicable.}

\def\TCInputData{%
  Errors set through the runtime (number 11, a description, a source, a help
  context), \texttt{Err.\allowbreak{}Raise} with and without a description,
  source and help file, \texttt{Err.\allowbreak{}Raise 0}, and
  \texttt{Err.\allowbreak{}Clear}.}
\def\TCInitialState{%
  No state from earlier runs is needed; the test makes what it needs.}

\def\TCProcedure{%
  \begin{enumerate}
    \item CTest runs \texttt{vbrun\_\allowbreak{}errobj\_\allowbreak{}tests}.
    \item Read \texttt{Number}, \texttt{Description}, \texttt{Source},
      \texttt{HelpContext} and \texttt{LastDllError} directly, raise and
      clear errors, and check that a new error clears the source and that
      \texttt{Err.\allowbreak{}Raise 0} is error 5.
    \item Ask the object for \texttt{IUnknown} and another interface, and look
      up the member names.
    \item Read the number and raise an error through late binding.
    \item CTest records the exit code and, on failure, the program's output
      (\texttt{--output-on-failure}).
  \end{enumerate}}
\def\TCExpectedResult{%
  All checks pass; the program prints \texttt{errobj ok} and exits 0.}
\def\TCPostConditions{%
  No file or registry state remains that later tests depend on.}

\TCTemplate

\end{document}
